\documentclass[10pt]{article}

\usepackage[utf8]{inputenc}

\usepackage[T1]{fontenc}

\usepackage{amsmath}

\usepackage{amsfonts}

\usepackage{amssymb}

\usepackage[version=4]{mhchem}

\usepackage{stmaryrd}

\usepackage{graphicx}

\usepackage[export]{adjustbox}

\graphicspath{ {./images/} }



\begin{document}

тинуум, т. е. связное компактное метрич. пространство $C$, каждая точка к-рого обладает сколь угодно малой окрестностью с границей размерности нуль. Другими словами, при любом $\varepsilon>0$ пространство $C$ может быть представлено в виде суммы конечного числа замкнутых множеств диаметра, меньшего $\varepsilon$, обладающих тем свойством, что никакие три из этих множеств не имеют общей точки. Ковер Серпиньского удовлетворяет әтому определению Л., так что всякая канторова кривая является также и Л. в смысле П. С. Урысона. Обратно, если плоский континуум является Л. в смысле П. С. Урысона, то он будет канторовой кри-



\begin{center}

\includegraphics[max width=\textwidth]{2023_07_08_13bdd5a9f6d51ce398a8g-1(3)}

\end{center}



Рис. 3. вой. Определение Л., данное П. С. Урысоном, является внутренним: оно характеризуется лишь свойствами самого пространства $C$ и не зависит от того, рассматривается ли это пространство само по себе или как подмножество другого топологич. пространства.



Существуют Л., к-рые не гомеоморфны никакому подмножеству плоскости. Такова, напр., Л., лежащая в трехмерном пространстве и состоящая из шести ребер тетраэдра и четырех отрезков, соединяющих какую-либо точку пространства, не лежащую ни на одной из его граней,



\begin{center}

\includegraphics[max width=\textwidth]{2023_07_08_13bdd5a9f6d51ce398a8g-1(2)}

\end{center}



Рис. 4. с его вершинами (рис. 3). Но всякая Л. (в смысле П. С. Урысона) гомеоморфна нек-рому подмножеству трехмерного евклидова пространства ( т е о р е м а М е н г е р а). Континуум $M$, обладающий тем свойством, что какова бы ни была Л. $C$, в $M$ найдется подконтинуум $C^{\prime}$, гомеоморфный континууму $C$, строится следующим образом. Куб $K$ с ребром 1 делится плоскостями, параллельными его граням, на 27 равных кубов. Из куба $K$ удаляются центральный куб и все прилежащие к нему по двумерным граням кубы этого подразделения. Получается множество $K_{1}$, состоящее из 20 оставшихся замкнутых кубов первого ранга. Поступая точно так же с каждым из кубов первого ранга, получим множество $K_{2}$, состоящее из 400 кубов второго ранга (рис. 4). Продолжая этот процесс бесконечно, получим последовательность континуумов $K_{1} \supset K_{2} \supset K_{3} \supset$. ., пересечение к-рых есть одномерный континуум $M$, наз. у н и в е р с а л ь н о й к р и в о й M е н г e p a.



В исследовании Л. важную роль играет понятие и нд е к с а в е т в ле ни я. Л. $C$ в точке $x$ имеет индекс ветвления $\mathfrak{m}$, если каково бы ни было число $\varepsilon>0$, существует открытое множество $U$ диаметра, меньшего, чем $\varepsilon$, содержащее точку $x$, граница к-рого есть множество мощности, не превосходящей $\mathfrak{m}$, но для достаточно малого $\varepsilon^{\prime}>0$ граница всякого открытого множества, содержащего точку $x$, диаметр к-рого меньше $\varepsilon^{\prime}$, имеет мощность, не меньшую, чем $\mathfrak{m}$. Точки Л. относительно их индекса ветвления классифицируются следующим образом.



\begin{enumerate}

  \item Точки с индексом ветвления $n, n$ - натуральное.



  \item Точки неограниченного индекса ветвления $\omega$. (Точка $x$ Л. $C$ имеет индекс ветвления $\omega$, если каково бы ни было число $\varepsilon>0$, существует открытое множество, содержащее точку $x$, с диаметром, меньшим, чем $\varepsilon$, граница к-рого состоит из конечного множества точек; но каково бы ни было натуральное $n$, найдется такое $\varepsilon_{n}>0$, что граница всякого открытого множества, содержащего $x$ и имеющего диаметр меньший, чем $\varepsilon_{n}$, состоит не менее, чем из $n$ точек.)



  \item Точки счетного индекса ветвления §..



  \item Точки континуального индекса ветвления с.



\end{enumerate}



Точка Л. $C$, индекс ветвления к-рой больше двух, наз. т о ч к о й в е т в л е н и я; точка, индекс ветвления к-рой равен единице, наз. к о нце в о й т о чк 0 й.



П р и м е р ы. а) Отрезок во всех своих внутренних точках имеет индекс ветвления, равный двум; индекс ветвления концов отрезка равен единице. б) Окружность в каждой своей точке имеет индекс ветвления два. в) Л., состоящая из $n$ прямолинейных отрезков,



\begin{center}

\includegraphics[max width=\textwidth]{2023_07_08_13bdd5a9f6d51ce398a8g-1}

\end{center}



Рис. 5 .



\begin{center}

\includegraphics[max width=\textwidth]{2023_07_08_13bdd5a9f6d51ce398a8g-1(1)}

\end{center}



Рис. 6 . исходящих из одной точки $O$, имеет в точке $O$ индекс ветвления n. г) Л., состоящая из отрезков $O a_{1}, O a_{2}$, . .., $O a_{n}, \ldots$. выходящих из начала координат $O$, имеюних длины $\frac{1}{2}, \frac{1}{2^{2}}, \ldots, \frac{1}{2^{n}}, \ldots$ и образующих с осью $O x$ углы, соответственно равные $\frac{\pi}{2}, \frac{\pi}{2^{2}}, \ldots$, $\frac{\pi}{2^{n}}$, . . имеет в точке $O$ неограниченно возрастающий индекс ветвления $\omega$ (рис. 5). д) Л., состоящая из отрезка $O a_{0}$ длины 1 и отрезков $O a_{1}, O a_{2}, \ldots, O a_{n}, \ldots$ длины 1 , выходящих из точки $O$ и образующих с отрезком $O a_{0}$ углы, соответственно равные $\frac{\pi}{2}, \frac{\pi}{2^{2}}, \ldots$, $\frac{\pi}{2^{n}}, \ldots$, имеет в каждой точке отрезка $O a_{0}$ счетный индекс ветвления ㄲo (рис. 6). е) Л., состоящая из отрезков, соединяющих точку $O$ со всеми точками канторова множества, лежащего на отрезке $0 \leqslant x \leqslant 1$,



\begin{center}

\includegraphics[max width=\textwidth]{2023_07_08_13bdd5a9f6d51ce398a8g-1(4)}

\end{center}



Puc. 7.



$y=0$, имеет во всех своих точках континуальный индекс ветвления $\mathfrak{c}$ (рис. 7). ж) Ковер Серпиньского также имеет во всех своих точках континуальный индекс ветвления.



Если у Л. совсем нет точек ветвления, т. е. если в каждой точке Л. индекс ветвления равен 1 или 2, то эта Л. есть либо п р о с т а я д у г а - топологич. образ отрезка, либо пр о с т а я 3 а м к н у т а я л и н и я - топологич. образ окружности. При этом, если индекс ветвления Л. во всех точках равен 2 , то это - простая замкнутая Л., если же у Л., не имеющей



$\triangle 13$ Математическая энц., т. 3





\end{document}